\section{B. The equations, their graded forms, and finite tests}

\subsection{The exact \S46 equations and the conversion to Drinfeld coordinates}

\begin{claim}{B.m1}{Proved}{the loop equation, with its generation condition}
Let $\widehat B=\widehat{C_2*C_3}$, with generators $s=\sigma_\infty$, $r=\rho$, and let $\widehat F$ be its level-two kernel. Put $e=sr$, $y=e^2=\ell_0$, $x=(rs)^2=\ell_1$. For $\mu\in\widehat{\mathbb Z}^{\times}$ put $\nu=(\mu-1)/2\in\widehat{\mathbb Z}$. A triple $h,g,f_G\in\widehat F$ specifies a loop automorphism inducing the identity on $S_3$ by
\[
 u(s)=hs,\quad u(r)=gr,\quad u(e)=f_G^{-1}e^\mu f_G
\]
if and only if
\begin{align*}
 h\sigma(h)&=1,&g\rho(g)\rho^2(g)&=1,\\
 \sigma(h)g&=\sigma(f_G)^{-1}x^\nu\rho(f_G),
\end{align*}
and $hs,gr$ topologically generate $\widehat B$. Here $f_G$ is defined modulo left multiplication by $y^n$, $n\in\widehat{\mathbb Z}$. The generation condition is part of the theorem; it is not a consequence claimed from the displayed identities.

Equivalently, after composing with $\operatorname{Int}(f_G)$, one obtains $u_0(s)=\eta s$, $u_0(r)=\xi r$ with
\begin{equation}\label{eq:LM46}
 \eta\sigma(\eta)=1,\qquad
 \xi\rho(\xi)\rho^2(\xi)=1,\qquad
 \xi=\eta x^\nu.
\end{equation}
The gauge transformation is
\[
 (\xi,\eta,f_G)\longmapsto
 (y^n\xi x^{-n},y^n\eta x^{-n},y^nf_G).
\]
These are \S46, fols.392--394 and 407--408, scans186--190 and 215--217, equations (25), (29), (44), (46), (133), (134), (144).

\end{claim}
\begin{proof}
The universal property of a free profinite product makes $s\mapsto hs$, $r\mapsto gr$ an endomorphism precisely when $(hs)^2=(gr)^3=1$, which expands into the two cocycles. It induces the identity modulo $\widehat F$. Its image is the closed subgroup generated by its two indicated images. A surjective endomorphism of a finitely generated profinite group is injective: for each $n$ there are only finitely many open normal subgroups of index at most $n$, and the inverse-image map permutes this finite set; an element in the kernel lies in all of them. This proves the role of the generation condition.

Expanding $u(sr)=f_G^{-1}(sr)^\mu f_G$ gives
$f_Gh\sigma(g\rho(f_G)^{-1})=y^\nu$.
Apply $\sigma$ and use $\sigma(y)=x$ to obtain the asserted equation. Set
$\eta=f_Gh\sigma(f_G)^{-1}$ and
$\xi=f_Gg\rho(f_G)^{-1}$.
The two cocycles telescope under this change, and the loop equation becomes $\xi=\eta x^\nu$. Conversely these substitutions recover $g,h$ and all the original equations. Finally $\rho(y)=\sigma(y)=x$ proves the gauge formula. Centralizers of the free generators are their procyclic closures, so the freedom in the conjugator for $y^\mu$ is exactly the stated $y^n$.
\end{proof}

\begin{claim}{B.m2}{Proved}{An explicit equivalence with (I) and (II)}
Modulo the gauge above, solutions of \eqref{eq:LM46} satisfying the generation condition are in bijection with the pairs $(\lambda,f)$ defining $\widehat{\mathrm{GT}}_0$, with $\lambda=\mu$ and
\[
 f\in\widehat F',\quad
 f(y,x)f(x,y)=1,\quad
 f(z,x)z^\nu f(y,z)y^\nu f(x,y)x^\nu=1.
\]
The pair acts by $x\mapsto x^\mu$, $y\mapsto f^{-1}y^\mu f$ and is required to be an automorphism. If the abelianization of $\xi$ is $a[x]+b[y]$, the exact change of variables is
\begin{equation}\label{eq:normalization-GT}
 a+b=\nu,\qquad f=y^{-b}\xi x^{-a}.
\end{equation}
An inverse is $\xi=y^\nu f$, $\eta=y^\nu fx^{-\nu}$.

\end{claim}
% anchor:B.m2
\manuscriptanchor{\S 46; fols393--394,407--408; scans189--190,215--217}

\begin{proof}
Abelianize $\eta=\xi x^{-\nu}$ and the $\sigma$-cocycle, where $\sigma$ interchanges $x,y$. Its sum is $(a+b-\nu)([x]+[y])=0$, so $a+b=\nu$. Thus \eqref{eq:normalization-GT} belongs to the closed commutator subgroup. Under gauge by $y^n$, the exponents become $(a-n,b+n)$ and the expression for $f$ stays unchanged. Choose $n=a$, giving $\xi=y^\nu f$, as claimed.

Substitution into the first equation gives
\[
 \eta\sigma(\eta)=y^\nu f(x,y)f(y,x)y^{-\nu},
\]
so it is equivalent to (I). Since $\rho(x,y,z)=(z,x,y)$, substitution into the second gives
\[
 y^\nu f(x,y)x^\nu f(z,x)z^\nu f(y,z)=1.
\]
Moving its first two factors cyclically to the end gives exactly (II). A cyclic permutation of a product is conjugate to that product; thus this step is an equivalence, not an implication after abelianization.

The automorphism $u_0$ fixes $y$ to $y^\mu$. Since $x=\rho(y)$, it sends $x$ to $\xi x^\mu\xi^{-1}$. Consequently $\operatorname{Int}(\xi^{-1})u_0$ restricts to
$x\mapsto x^\mu$, $y\mapsto\xi^{-1}y^\mu\xi=f^{-1}y^\mu f$.
Conversely a pair satisfying (I),(II) makes \eqref{eq:LM46} define an endomorphism of $\widehat B$. Its restriction to $\widehat F$ is conjugate to the stated automorphism and its quotient on $S_3$ is the identity; it is therefore surjective and hence an automorphism. This proves both directions, including the generation condition. The change of variables in arbitrary gauge is implemented on genuine automorphisms by $\operatorname{Int}(x^a\xi^{-1}f_G)u$.
\end{proof}

This proof is independent of any proposed graded identification. It agrees with Harbater--Schneps (2000), \S1.2 Proposition 2 and \S1.3 Propositions 3--4, Theorem 2, which we read. Their normalized lift supplies the same two-cycle and three-cycle relations. There is no five-point configuration in the proof.

\begin{claim}{B.m3}{Proved}{the additional double-marked condition has separate content}
The \S47 condition is compatibility with the specific continuous surjection
\[
 \varphi:\widehat F\twoheadrightarrow\widehat B,
 \qquad (y,x,z)\longmapsto(sr,s,r^{-1}).
\]
For a representative $u_0$ fixing $e$ to $e^\mu$, strict commutation $\varphi\circ(u_0|_{\widehat F})=u_0\circ\varphi$ is equivalent to
\begin{equation}\label{eq:additional-phi}
 \xi=\bigl(\varphi(\xi)\sigma(\varphi(\xi))^{-1}\bigr)x^\nu.
\end{equation}
In particular, deriving (I),(II) from \S46 does not prove that every element of $\widehat{\mathrm{GT}}_0$ satisfies \S47. This is the precise extra condition defining $C^\varphi$ in the dictionary. The passage is \S47, fols.416--420, scans231--240, especially (14),(15),(20),(30).

\end{claim}
\begin{proof}
The elements $x,y$ freely generate $\widehat F$, so equality of the homomorphisms need only be checked there. For $y$, the two images are $e^\mu$. For $x$, the two images are
$\varphi(\xi)s^\mu\varphi(\xi)^{-1}$ and $\eta s$.
Since $s$ has order two and $\mu$ is odd, equality is exactly
$\eta=\varphi(\xi)\sigma(\varphi(\xi))^{-1}$.
Use $\xi=\eta x^\nu$ to get \eqref{eq:additional-phi}. The subgroup of strictly commuting representatives is closed under composition and inversion because this is an equality of typed homomorphisms. Outer commutation allows a further conjugating element in $\widehat B$; it is a different assertion until a normalization has been made.
\end{proof}

\begin{claim}{B.m4}{Proved}{Proved in the stated representatives; global classification open}
The three fundamental cases of \S46 are: $u(\sigma)=\sigma$, equivalent to $h=1$ and, when a coboundary conjugator is chosen, $\beta=f_G$; $u(\rho)=\rho$, equivalent to $g=1$ and $\alpha=f_G$; and normalization of $\langle\ell_0\rangle$, equivalent to taking $f_G=1$ after the allowed change of representative. If $K_{1/2}$ is the centralizer section, then
\[
 K_2=\rho K_{1/2}\rho^{-1},\qquad
 K_{-1}=\rho^{-1}K_{1/2}\rho
\]
with the same formulas for the homomorphisms after identifying their common outer quotient. The order-three centralizer can project only to $\mu\equiv1\pmod3$.

\end{claim}
% anchor:B.m4
\manuscriptanchor{\S 46; fols402--405 verso; scans206--213}

\begin{proof}
The equivalences follow by substituting the definitions of $g,h,\alpha_0,\beta_0$. The cyclic action on the three fixed points is the same action that permutes the three puncture labels, so conjugating the subgroup and its section gives the two displayed sections; no choice of a new path is made. The abelianization of $\widehat B$ is $C_2\times C_3=C_6$. A loop automorphism sends the image of $e$ to its $\mu$th power. Its $C_3$ coordinate is the coordinate of $r$. If $u(r)=r$, that coordinate is unchanged, forcing $\mu\equiv1\pmod3$. Conversely, the required adjustment to a centralizing representative is a nonabelian $H^1$ problem, treated with its ambiguity in Part D; a matrix cocycle alone does not prove that adjustment.
\end{proof}

\begin{claim}{B.m5}{Refuted}{A conjugator in the smaller group does not define an outer condition}
The literal condition on \S47, fol.415, scan229 (4), that the conjugator $\Lambda$ for
$\varphi u=\operatorname{Int}(\Lambda)u\varphi$ belong to $\widehat F$, is not invariant under changing a representative by $\operatorname{Inn}(\widehat F)$. It cannot, with that literal target, define the claimed inverse image of an outer subgroup containing every geometric inner automorphism. The correct ordinary outer-compatibility condition allows $\Lambda\in\widehat B$. The stronger condition $\Lambda\in\widehat F$ is a separate relative condition. The margin on scan230 already notices that $\varphi(f)$ need not lie in $\widehat F$.
\end{claim}
\begin{proof}
Take $u=\operatorname{Int}(x)$, where $x=\ell_1\in\widehat F$ and $\varphi(x)=s$. Then
$\varphi u=\operatorname{Int}(s)\varphi$ and
$u\varphi=\operatorname{Int}(x)\varphi$.
Surjectivity of $\varphi$ and $Z(\widehat B)=1$ imply that the only possible conjugator is $\Lambda=sx^{-1}$. Its reduction in $B/F=S_3$ is the nontrivial transposition $s$, so it is not in $\widehat F$. The identity representative of the same outer class has conjugator $1\in\widehat F$. This is a complete counterexample to representative invariance; the program checks the obstructing reduction modulo two. Allowing $\Lambda\in\widehat B$ removes precisely this obstruction, since inner changes on either side merely change a target inner conjugator. Statement B.m3 uses strict commutation for an explicitly normalized representative, and remains valid without either ambiguity.
\end{proof}

\begin{claim}{B.m6}{Proved}{The alternate \S47 formulas and their finite labels}
Assume the system of B.m1, the generation condition, and strict $\varphi$-compatibility for $u_0$. Put $a_3=(1-\varepsilon_3(\mu))/2$ and $\widetilde\alpha=\alpha J^{a_3}\in\widehat F$. There are unique $i\in\mathbb Z/2$ and $j'\in\mathbb Z/3$ such that
\[
 \varphi(\xi)=\beta s^i,\qquad
 \varphi\bigl(\rho^{-1}(\xi)\bigr)=r^{j'}s^{a_3}\widetilde\alpha^{-1}.
\]
Conversely either displayed compatibility formula, together with the loop system and its generation condition, implies the other. Reduction in $B/F=S_3$ gives exactly (51) of scan247 and uniquely determines $i,a_3,j'$ from these two reductions. Moreover $i\equiv\nu\pmod2$, so $i=(1-\varepsilon_2(\mu))/2$; no analogous determination of $j'$ from $\mu$ is asserted here.

When $\mu\equiv1\pmod3$, the full compatible loop system can equivalently be written as the \emph{conjunction}
\begin{align*}
 \xi&=w\rho(w)^{-1},\qquad w=\varphi\bigl(\rho^{-1}(\xi^{-1})\bigr),\\
 \xi&=\varphi(\xi)\sigma(\varphi(\xi))^{-1}x^\nu,
\end{align*}
again retaining the generation condition. These are (57),(58), \S47, fols.423--425, scans245--250. The reading ``l'ensemble'' on scan250 is supported by this equivalence. We do not replace it by ``l'une'' without specifying additional standing hypotheses.
\end{claim}
\begin{proof}
Statement B.m3 says that checking strict compatibility on the free generator $x$ is equivalent to
$\varphi(\xi)s\varphi(\xi)^{-1}=\beta s\beta^{-1}$. The centralizer of $s$ in $\widehat B$ is $\langle s\rangle$, so this is equivalent to the first displayed formula, with unique $i$.

One may instead check on the free basis $(y,z)$. The $y$ equation holds by normalization. The formula $z=r^{-1}yr$ gives
$u_0(z)=\operatorname{Int}(\rho^{-1}(\xi^{-1}))(z^\mu)$. Write $v=\varphi(\rho^{-1}(\xi))$. Compatibility on $z$, using $\varphi(z)=r^{-1}$, is
$v\alpha r\alpha^{-1}v^{-1}=r^\mu=r^{(-1)^{a_3}}$.
The element $\kappa=sJ$ inverts $r$. Hence $\kappa^{a_3}v\alpha$ centralizes $r$, so equals $r^j$ for a unique $j\in\mathbb Z/3$, by the free-product centralizer theorem used in C.d1. Solving and setting $j'=(-1)^{a_3}j$ yields
$v=r^{j'}s^{a_3}(\alpha J^{a_3})^{-1}$, as claimed. This proves the equivalence by checking the two free bases, with all comparison maps specified. Reduction gives the stated finite formulas because $\widetilde\alpha,\beta$ belong to $\widehat F$. To determine $i$, the sign character of $S_3$ sends both $\varphi(x)$ and $\varphi(y)$ to $-1$. Thus $(-1)^i=\operatorname{sgn}(\varphi(\xi))=(-1)^{a+b}=(-1)^\nu$, by B.m2. In the unnormalized formula $\varphi(\xi)=e^r\beta s^i$, the same argument gives $i+r\equiv\nu\pmod2$; setting $r=0$ by gauge gives the stated value.

If $a_3=0$, the second formula says $w=\alpha r^{-j'}$; taking its $\rho$-coboundary gives the first of (57),(58). The second is B.m3. Conversely the first is a coboundary and therefore satisfies the $\rho$-cocycle by telescoping. The second makes $\eta=\varphi(\xi)\sigma(\varphi(\xi))^{-1}$ a $\sigma$-coboundary and supplies the loop relation. B.m1 with the generation condition reconstructs the automorphism, and B.m3 proves strict compatibility. Thus the conjunction is exactly sufficient as well as necessary.
\end{proof}

\begin{claim}{B.m7}{Known}{Complex conjugation is self-centralizing already in $\widehat{\mathrm{GT}}_0$}
Let $c=(-1,1)$ in the normalized pair group. Then
$C_{\widehat{\mathrm{GT}}_0}(c)=\{1,c\}$.
The profinite fixed-subgroup input is
$\widehat F^{c}=1$ for $c(x)=x^{-1},c(y)=y^{-1}$, Lochak--Schneps (1997), Proposition 4(i); their Proposition 4(ii) proves self-centralization in $\widehat{\mathrm{GT}}$ \citep{LochakSchneps1997dictionary}. We opened both statements and their proof in \S2, printed p.581; the latter uses relation (II) and no pentagon. We derive the larger-group assertion using the genuine loop equations, without the pentagon. This proves the proposition of \S47, fols.426--427, scans252--253, and corrects its final $\nu=\pm1$ to $\nu\in\{0,-1\}$.
\end{claim}
\begin{proof}
One can identify the fixed-subgroup input with the finite-factor centralizer theorem: the semidirect product $\widehat F\rtimes\langle c\rangle$ has the free profinite product presentation $C_2\amalg C_2\amalg C_2$, on $c,xc,yc$. The centralizer of the first factor's nontrivial element is that factor; its intersection with $\widehat F$ is trivial. This is the same standard profinite centralizer input isolated in D.3; Proposition 4(i) of the opened primary article gives the needed consequence directly.

For a normalized pair $(\mu,f)$ commuting with $c$, compare their genuine normalized representatives. On $y$ this gives conjugation of $y^{-\mu}$ by $c(f)$ and by $f$. Their quotient centralizes $y$, so is a power of $y$; its abelianization is zero because $f,c(f)\in\widehat F'$. Thus $c(f)=f$, and the fixed-subgroup result forces $f=1$. The inverse construction in B.m2 first gives $\xi=y^\nu$, $\eta=y^\nu x^{-\nu}$. Gauge by $y^{-\nu}$ changes these to $\xi=x^\nu$, $\eta=1$, without changing the pair.

For this actual profinite solution, C.d1 supplies the matrix lifts $\eta=I$ and $\xi=\varepsilon L_1^\nu$ with one global sign $\varepsilon$. Its trace constraint is
$1=\operatorname{tr}(\xi R)=\varepsilon(2\nu+1)=\varepsilon\mu$.
Hence $\mu=\varepsilon$ is globally $1$ or $-1$, not an independently chosen sign at each prime. These two pairs are precisely $1,c$, and both commute with $c$. This completes the reduction and the claimed correction of the manuscript's last line.
\end{proof}

